\clearpage
\VerdierPhysicalPageAnchor{0235}%
\begingroup\small\setlength{\parskip}{0pt}\setlength{\emergencystretch}{1.5em}
\setlength{\abovedisplayskip}{4pt plus 2pt minus 2pt}%
\setlength{\belowdisplayskip}{4pt plus 2pt minus 2pt}%

\VerdierSemanticAnchor{ch4d-4-2-2-cont}
\noindent où $(p,q)$ et $(p',q')$ sont des éléments de
$\mathcal{F}\ell_1(\widetilde{\mathbb Z})$ tels que $p\leq p'$ et
$q\leq q'$, l'image dans $D(\mathcal A)$ du morphisme canonique de
$\operatorname{comp}(\mathcal A)$ :

\VerdierSemanticAnchor{ch4d-4-2-2-3}
\begin{equation}
\psi\bigl((p,q),(p',q')\bigr)\colon
X(p,q)\longrightarrow X(p',q').
\tag{4.2.2.3}
\end{equation}
On a ainsi déterminé un foncteur de
$\mathcal{F}\ell_1(\widetilde{\mathbb Z})$ dans $D(\mathcal A)$. Soit,
pour tout élément $(p,q,r)$ de
$\mathcal{F}\ell_2(\widetilde{\mathbb Z})$ :

\VerdierSemanticAnchor{ch4d-4-2-2-4}
\begin{equation}
S(p,q,r)\colon
0\longrightarrow X(p,q)\longrightarrow X(p,r)
\longrightarrow X(q,r)\longrightarrow0,
\tag{4.2.2.4}
\end{equation}
la suite exacte de $\operatorname{comp}(\mathcal A)$ dont les morphismes
sont les morphismes canoniques déduits de la filtration de $X$. Soit enfin,
pour tout élément $(p,q)$ de
$\mathcal{F}\ell_1(\widetilde{\mathbb Z})$,

\VerdierSemanticAnchor{ch4d-4-2-2-5}
\begin{equation}
J(p,q)\colon\mathcal X(p,q)\longrightarrow X(p,q),
\tag{4.2.2.5}
\end{equation}
l'image dans $D(\mathcal A)$ du morphisme de
$\operatorname{comp}(\mathcal A)$ :

\VerdierSemanticAnchor{ch4d-4-2-2-6}
\begin{equation}
j(p,q)\colon\mathcal X(p,q)\longrightarrow X(p,q),
\tag{4.2.2.6}
\end{equation}
dont les composants sont :

\VerdierSemanticAnchor{ch4d-4-2-2-7}
\begin{equation}
\left\{
\begin{aligned}
j(p,q)^i\colon\mathcal X(p,q)^i
 &=X(q)^i\oplus X(p)^{i+1}
   \longrightarrow X(q)^i/X(p)^i=X(p,q)^i,\\
j(p,q)^i&=(\pi^i,0),\\
&\text{où $\pi^i$ est l'épimorphisme canonique :}\\
\pi^i\colon X(q)^i&\longrightarrow X(q)^i/X(p)^i.
\end{aligned}
\right.
\tag{4.2.2.7}
\end{equation}

\VerdierSemanticAnchor{ch4d-4-2-3}
\medskip\noindent\textbf{Proposition 4.2.3.}\quad
\textit{Pour tout élément $(p,q,r)$ de
$\mathcal{F}\ell_2(\widetilde{\mathbb Z})$, désignons par :}
\[
\delta(p,q,r)\colon X(q,r)\longrightarrow X(p,q)[1],
\]
\noindent\textit{le morphisme de $D(\mathcal A)$ :}

\VerdierSemanticAnchor{ch4d-4-2-3-1}
\begin{equation}
\delta(p,q,r)=\delta\bigl(S(p,q,r)\bigr),
\tag{4.2.3.1}
\end{equation}
\noindent\textit{où $\delta\bigl(S(p,q,r)\bigr)$ est le morphisme de
$D(\mathcal A)$ associé à la suite exacte $S(p,q,r)$ (4.2.2.4) par le
$\delta$-foncteur canonique
$\operatorname{comp}(\mathcal A)\longrightarrow D(\mathcal A)$ (1.3.2).}

\VerdierSemanticAnchor{ch4d-4-2-3-ab}
\noindent\textit{a) Les objets $X(p,q)$ (4.2.2.1), les morphismes
$\varphi\bigl((p,q),(p',q')\bigr)$ (4.2.2.2) et les morphismes
$\delta(p,q,r)$ (4.2.3.1) définissent un objet spectral à valeurs dans
$D(\mathcal A)$.}

\noindent\textit{b) Les morphismes $J(p,q)$ (4.2.2.5) définissent un
isomorphisme de l'objet spectral
$\bigl(\mathcal X(p,q),\delta(p,q,r)\bigr)$ sur l'objet spectral
$\bigl(X(p,q),\delta(p,q,r)\bigr)$.}

\VerdierSemanticAnchor{ch4d-4-2-3-proof}
L'assertion a) résulte immédiatement des propriétés du $\delta$-foncteur
canonique
$\operatorname{comp}(\mathcal A)\longrightarrow D(\mathcal A)$ (1.3.1).
Pour démontrer l'assertion b), le lecteur se reportera à (1.3.2).
\endgroup
